% Stable corpus identity 2590; LMCS/v21i3/16491/Theorem-3.2.
% Faithful editable transcription of the complete selected author proof and context.
% Source PDF SHA256: 622f85d1ec0aab4ee76e79197f37a3f2e26aba7b72e7b6e7da2550b896553994.
% Source order of authored excerpts: 26:3, 26:7, 26:8, 26:9, 26:10, 26:16.
% No mathematical repair, rounding convention, new proof, or omitted proof core supplied.
\documentclass[11pt]{article}
\usepackage[a4paper,margin=25mm]{geometry}
\usepackage{fontspec}
\setmainfont{DejaVu Serif}
\usepackage{amsmath,amssymb,amsthm}
\usepackage[hidelinks]{hyperref}
\usepackage{enumitem,needspace}
\setlength{\parindent}{1.3em}
\setlength{\parskip}{3pt}
\setlength{\emergencystretch}{2em}
\newcommand{\NP}{\mathrm{NP}}
\newcommand{\coNP}{\mathrm{coNP}}
\newcommand{\Pclass}{\mathrm{P}}
\newcommand{\SAT}{\mathrm{SAT}}
\newcommand{\SIZE}{\mathrm{SIZE}}
\newcommand{\NSIZE}{\mathrm{NSIZE}}
\newcommand{\coNSIZE}{\mathrm{coNSIZE}}
\newcommand{\NTIME}{\mathrm{NTIME}}
\newcommand{\poly}{\mathrm{poly}}
\newcommand{\size}{\mathrm{size}}
\newcommand{\sourcepage}[1]{\par\Needspace{5\baselineskip}\bigskip\noindent\textit{[Original source page 26:#1]}\par\medskip}
\newtheorem*{sourcetheoremthreetwo}{Theorem 3.2 \textnormal{(Restatement of Theorem 1.2 Item $(i)$)}}
\newtheorem*{sourceclaimthreethree}{Claim 3.3}
\begin{document}
\section*{Editorial record: problem 2590}
\noindent\textbf{Selected problem.} For rational $1<k<c<k^2$, there is no eventually successful non-uniform refuter of circuit size at most $n^c$ for $\SAT_n$ against co-nondeterministic circuits of size $n^k$, exactly as stated in Theorem 3.2 below. Claim 3.3 and all context remain uncounted.

\noindent\textbf{Difficulty.} A refuter output need not reveal the truth of its SAT instance efficiently, so one use of the refuter does not yield the desired contradiction. The author's proof recursively packages the complement of a longer instance into a smaller-input circuit, applies refuters across shrinking input lengths, and resolves the final instance using a hardcoded truth table. It coordinates the opposing input-size and refuter-size constraints using $c<k^2$, controlling the stages and final table so that the resulting deciding circuit is itself small enough to be refuted.

\noindent\textbf{Attribution and version.} Lijie Chen, Jiatu Li, and Igor C. Oliveira, \emph{On the Unprovability of Circuit Size Bounds in Intuitionistic $\mathsf{S}^1_2$}, Logical Methods in Computer Science, volume 21, issue 3, article 26 (2025), pp.\ 26:1--26:20; submitted April 22, 2024; published September 10, 2025. DOI: \href{https://doi.org/10.46298/lmcs-21(3:26)2025}{10.46298/lmcs-21(3:26)2025}. \par\noindent Official article: \url{https://lmcs.episciences.org/16491}.\par\noindent Source PDF SHA256:\par\noindent\texttt{622f85d1ec0aab4ee76e79197f37a3f2e26aba7b72e7b6e7da2550b896553994}.

\noindent\textbf{Reuse and transcription.} The individual article is licensed under \href{https://creativecommons.org/licenses/by/4.0/}{Creative Commons Attribution 4.0 International}. The following original English excerpts reproduce the complete selected statement and human-authored proof, including the complete supporting Claim 3.3, and the exact SAT/refuter/circuit definitions in source order. Only unrelated source material is outside this packet. The author's low-order and polylogarithmic-overhead conventions, displayed parameters, and base case are preserved. Source-page markers and this opening record are editorial; no additional mathematical argument is supplied.

\clearpage
\section*{Original human-authored material}
\sourcepage{3}
Elaborating on our approach, we now present some new consequences for refuters that might be of independent interest. We use $\SAT_n$ to denote the decision version of the satisfiability problem for De Morgan Boolean formulas represented by $n$-bit strings (our results are robust to encoding details). As in the discussion above, we say that a refuter $R$ for a language $L$ succeeds against a class $\mathcal C$ of devices on input length $n$ if, for every $E\in\mathcal C$, $R(1^n,E)$ outputs $x\in\{0,1\}^n$ such that $E(x)\ne L(x)$. Observe that the existence of such a refuter implies that $L\notin\mathcal C$.

\sourcepage{7}
\noindent\textbf{2.2. Complexity theory.} We assume basic familiarity with complexity theory, e.g., complexity classes $\Pclass$, $\NP$ and $\coNP$ (see \cite{AB09}). We use $\NSIZE[s(n)]$ (resp.\ $\coNSIZE[s(n)]$) to denote the set of languages decidable by families of nondeterministic (resp.\ co-nondeterministic) circuits of size $s(n)$. For a set $\mathcal C$ of languages, we let $\mathrm{i.o.}\text{-}\mathcal C$ denote the set of languages $L'$ for which there is some $L\in\mathcal C$ such that $L'$ and $L$ agree on infinitely many input lengths.

\noindent\textbf{Uniform circuits.} Let $\mathcal C$ be a non-uniform complexity class, e.g., $\mathcal C=\SIZE[\poly(n)]$ or $\mathcal C=\NSIZE[\poly(n)]$. We say that $L\in\Pclass\text{-uniform }\mathcal C$ if $L$ is decidable by a family of $\mathcal C$-circuits $\{C_n:\{0,1\}^n\to\{0,1\}\}_{n\in\mathbb N}$ such that there is a polynomial-time Turing machine $M$ such that $M(1^n)$ outputs the description of $C_n$.

\noindent\textbf{Refuters.} Let $L$ be a language, $R$ be an algorithm, and $\mathcal C$ be a complexity class. We often abuse notation and also view $\mathcal C$ as a class of computational devices. We say that $R$ is a refuter for the lower bound $L\notin\mathcal C$ (or a refuter for $L$ against $\mathcal C$) on input length $n$ if, for every $E\in\mathcal C$, $R(1^n,E)$ outputs $x\in\{0,1\}^n$ such that $E(x)\ne L(x)$. We extend this definition in the natural way when $R$ represents a non-uniform family of circuits. In this case, we might omit the input $1^n$ and simply write $R_n$.

\sourcepage{8}
In the proof of the next result, we employ a more sophisticated \emph{bootstrapping argument} consisting of iterated applications of the refuter over different input lengths. Recall that we use $\SAT_n$ to denote the satisfiability problem for De Morgan Boolean formulas represented by $n$-bit strings.

\begin{sourcetheoremthreetwo}
Let $k$ and $c$ be rational numbers such that $1<k<c<k^2$. Then there is no non-uniform refuter $R(1^n,E_n)$ for $\SAT_n$ against $E_n\in\coNSIZE[n^k]$ that has circuit size $\le n^c$ and succeeds on every large enough input length $n$.
\end{sourcetheoremthreetwo}
\begin{proof}
In order to simplify some calculations, we prove the result for a fixed but arbitrary language $L\in\NTIME[n]$. Since $\poly(\log n)$ overheads in our estimates do not affect the result, it is easy to check that the same proof works for a language in $\NTIME[n\cdot\poly(\log n)]$, such as $\SAT=\{\SAT_n\}$. Similarly, we will tacitly assume that machines of complexity $t$ can be simulated by circuit of size $t$ (instead of $O(t\cdot\log t)$). Our construction and upper bounds can be easily adjusted to account for these small overheads, since for constants $a<b$, $n^a\cdot\poly(\log n)<n^b$ for every large enough $n$.

Let $1<k<c<k^2$, and let $M$ be a linear time nondeterministic machine that decides $L$. Now consider an arbitrary $n_0\in\mathbb N$, and suppose towards a contradiction that $R=\{R_\ell\}_{\ell\ge n_0}$ is a sequence of nonuniform circuits $R_\ell$ of size $\le\ell^c$ such that, given a circuit $E_\ell\in\coNSIZE[\ell^k]$

\sourcepage{9}
over $\ell$ input bits, $R_\ell(E_\ell)$ outputs a string $x_\ell$ such that $M(x_\ell)\ne E_\ell(x_\ell)$. We will prove the following claim.

\begin{sourceclaimthreethree}
There is an input length $\ell_0\ge n_0$ and a deterministic circuit $B_{\ell_0}$ of size at most $\ell_0^k$ that agrees with the language $L$ over $\{0,1\}^{\ell_0}$, i.e., $B_{\ell_0}(x)=M(x)$ for all $x\in\{0,1\}^{\ell_0}$.
\end{sourceclaimthreethree}
Note that the existence of $B_{\ell_0}$ is in contradiction with the existence of the refuter $R_{\ell_0}$, since there is no string $x_{\ell_0}$ such that $B_{\ell_0}(x_{\ell_0})\ne M(x_{\ell_0})$. Consequently, in order to complete the proof of Theorem 3.2, it is sufficient to establish Claim 3.3.

\begin{proof}[Proof of Claim 3.3]
We will consider a large enough input length $\ell_0\ge n_0$ specified later in the proof. First, we describe the pseudocode of a (recursively defined) deterministic circuit $B_{\ell_0}(x)$ that computes as follows:
\begin{enumerate}[label=(\arabic*)]
\item $B_{\ell_0}$ is given a string $x_{\ell_0}\in\{0,1\}^{\ell_0}$ as input.
\item It computes the description of a co-nondeterministic circuit $E_{\ell_1}(z)$ operating on inputs of length $\ell_1$ that ignores its input string $z$ and satisfies $E_{\ell_1}(z)\triangleq1-M(x_{\ell_0})$.

(The value $\ell_1<\ell_0$ will be defined below.)
\item $B_{\ell_0}$ obtains $x_{\ell_1}\triangleq R_{\ell_1}(E_{\ell_1})$.

(Note that if $\size(E_{\ell_1})\le\ell_1^k$ then $M(x_{\ell_1})=1-E_{\ell_1}(x_{\ell_1})=1-(1-M(x_{\ell_0}))=M(x_{\ell_0})$.)
\item Finally, it computes $b_{\ell_1}\triangleq M(x_{\ell_1})$, and uses this bit to decide $M(x_{\ell_0})$.

(As explained below, the key point of the construction is that we can (deterministically) compute the bit $b_{\ell_1}$ in a recursive fashion via Steps 1--3 using the sequence of refuters.)
\end{enumerate}
Before elaborating on the recursion performed in Step 4, we discuss the complexity parameters involved. Let $t(\ell_1)$ denote the circuit complexity of computing the bit $b_{\ell_1}\triangleq M(x_{\ell_1})$ on an arbitrary input $x_{\ell_1}$, i.e., the complexity of deciding if $x_{\ell_1}\in L$ over strings of length $\ell_1$. As explained below, in order for $B_{\ell_0}$ to be correct and of the desired size, it is sufficient that:
\begin{itemize}
\item $\size(E_{\ell_1})\approx\ell_0\le\ell_1^k$ (i.e., the refuter receives a circuit of bounded size in Step 3), where we have used that $M$ runs in nondeterministic linear time and omitted polylog factors.

(In our analysis below, we can assume that the description of $E_{\ell_1}$ can be computed in size at most $\ell_0^k/4$, since $k>1$ and $M$ runs in linear time.)
\item We need that $\ell_1\ge n_0$, so the output of $R_{\ell_1}(E_{\ell_1})$ is defined in Step 3.
\item The circuit size of the refuter in Step 3, which is upper bounded by $\ell_1^c$, satisfies $\ell_1^c\le\ell_0^k/4$.
\item The circuit size $t(\ell_1)$ needed to compute $b_{\ell_1}$ in Step 4 satisfies $t(\ell_1)\le\ell_0^k/4$.
\end{itemize}
If these conditions are met, $B_{\ell_0}(x)=M(x)$ for all $x\in\{0,1\}^{\ell_0}$, and its overall size $t(\ell_0)$ is strictly less than $\ell_0^k$ (with some room to spare that will be handy later in the proof), since
\[
t(\ell_0)\le\ell_0^k/4+\ell_0^k/4+t(\ell_1)\le\ell_0^k/4+\ell_0^k/4+\ell_0^k/4<\ell_0^k.
\]
\noindent\textbf{Constraints.} The conditions described above yield the following inequalities (omitting some low order terms):
\[
\ell_0^{1/k}\le\ell_1\le\ell_0^{k/c}\quad(\text{thus }c<k^2)\qquad\text{and}\qquad t(\ell_1)\le\ell_0^k/4.
\]
Recall that the condition $c<k^2$ is one of our assumptions. Jumping ahead, we will define $\ell_1$ as a function of $\ell_0$, $k$, and $c$, and employ a recursive approach to compute $b_{\ell_1}$ so that the conditions are satisfied.

\noindent\textbf{Key insight.} Note that in order to compute $b_{\ell_1}=M(x_{\ell_1})$ we must solve the \emph{same problem} on a \emph{smaller input length}, i.e., we would like to design a circuit $B_{\ell_1}$ that computes $L$ on

\sourcepage{10}
input length $\ell_1\ll\ell_0$. Consequently, using that we have refuter circuits for all input lengths $\ge n_0$, it is enough to iterate the same construction!

\noindent\textbf{Recursion.} The circuit $B_{\ell_1}$ computes analogously to $B_{\ell_0}$, by considering a smaller input length $\ell_2\ll\ell_1$ and a corresponding co-nondeterministic circuit $B_{\ell_2}$. Similarly to the analysis from above, we obtain the following constraints to guarantee its correctness and the desired circuit size bound (again omitting small order factors to focus on the relevant asymptotics):
\[
\size(E_{\ell_2})\approx\ell_1\le\ell_2^k\qquad\text{and}\qquad\ell_2^c\le\ell_1^k/4\qquad\text{and}\qquad t(\ell_2)\le\ell_1^k/4.
\]
This forces that
\[
\ell_1^{1/k}\le\ell_2\le\ell_1^{k/c}\qquad\text{and}\qquad
t(\ell_1)\le\ell_1^k/4+\ell_1^k/4+t(\ell_2)<3\cdot\ell_1^k/4\le\ell_0^k/4,
\]
where we assumed that $t(\ell_2)\le\ell_1^k/4$ and used that $\ell_1\ll\ell_0$. In other words, we need
\[
\ell_1^{1/k}\le\ell_2\le\ell_1^{k/c}\quad(\text{thus }c<k^2)\qquad\text{and}\qquad t(\ell_2)\le\ell_1^k/4.
\]
\noindent\textbf{Parameters and base case.} Consider the input lengths $\ell_0>\ell_1>\ldots>\ell_d$ explored in this way, together with the corresponding circuits $B_{\ell_0},\ldots,B_{\ell_d}$, where each $B_{\ell_i}$ contains $B_{\ell_{i+1}}$ as a subroutine. In other words, $B_{\ell_0}$ appears close to the input string, followed by $B_{\ell_1}$, and so on. (The string $x_{\ell_{i+1}}$ computed by $B_{\ell_i}$ serves as the input string to $B_{\ell_{i+1}}$.) We start with an input length $\ell_0$ sufficiently larger than $n_0$ so that all calls to the non-uniform refuter $R$ consider input lengths not smaller than $n_0$. In order to satisfy the inequalities from above, our parameters are defined as follows:
\begin{itemize}[label=--]
\item \emph{Sequence of input lengths.} We let $\ell_{i+1}\triangleq\ell_i^{1/k}$. Therefore, $\ell_d=\ell_0^{1/k^d}$.
\item \emph{Number of stages.} We take $d$ large enough, so that $\ell_d=\ell_0^{1/k^d}=\log\ell_0$, i.e., $d\triangleq(\log\log\ell_0-\log\log\ell_d)/\log k=O(\log\log\ell_0)$.
\item \emph{Initial input length $\ell_0$.} We want $\ell_d=\log\ell_0\ge n_0$, so we set $\ell_0\triangleq2^{n_0}$.
\end{itemize}
In the base case (input length $\ell_d$ and circuit $B_{\ell_d}$), we simply consult the hardcoded truthtable of the language $L$ computed by machine $M$ on inputs of length $\ell_d=\log\ell_0$. The truthtable of $L$ on input length $\log\ell_0$ can be nonuniformly stored using $\ell_0$ bits. This can be seen as an overhead in the final size of $B_{\ell_0}$ that is of order $\ell_0\le\ell_0^k/4$, where this inequality uses that $k>1$ and assumes that $\ell_0\ge n_0$ is large enough. Since the overall size bound for $B_{\ell_0}$ is given by a simple additive function obtained from the concatenation of a sequence of circuits, it is not hard to see that its total size is at most $\ell_0^k$, as desired.
\end{proof}
As explained above, this completes the proof of Theorem 3.2.
\end{proof}

\sourcepage{16}
\begin{thebibliography}{AB09}
\bibitem[AB09]{AB09} Sanjeev Arora and Boaz Barak. \emph{Computational Complexity - A Modern Approach}. Cambridge University Press, 2009. URL: \url{http://www.cambridge.org/catalogue/catalogue.asp?isbn=9780521424264}.
\end{thebibliography}
\end{document}
